We compare the \texttt{bisect} cost estimate against two benchmarks: the Iowa Legislative Services Agency model and the California Citizens Redistricting Commission.

\subsection{The Iowa Model: \$3 Million Per Cycle}

Iowa's Legislative Services Agency is widely regarded as the most cost-effective non-partisan redistricting model in the United States. Its 2021 redistricting cycle cost approximately \$3 million, including:

\begin{itemize}
  \item LSA staff time (three full-time staff over 8 months): approximately \$480,000
  \item Outside counsel (constitutional and VRA review): approximately \$200,000
  \item Public hearings (statewide; multiple rounds of maps): approximately \$150,000
  \item Legislative review process (committee hearings, floor debate): approximately \$400,000
  \item Technical systems (GIS software, mapping tools): approximately \$75,000
  \item Document publication and public engagement: approximately \$50,000
  \item Expert demographic analysis: approximately \$120,000
  \item Contingency and administrative overhead: approximately \$1,525,000
\end{itemize}

The \$3 million Iowa estimate includes the cost of Iowa's three-map process (the LSA draws and redraws if rejected) and the legislative review process. The dominant cost is staff time for a process that takes 8 months, compared to the \texttt{bisect} implementation's 12 weeks.

The algorithmic approach eliminates most of the map-drawing staff time (the algorithm draws the baseline in minutes, not months) and reduces the number of maps drawn (one baseline plus documented modifications, not multiple redrafts). The legislative review process cost is unaffected (it is a function of legislative procedures, not algorithm design), but is computable only if the state retains a legislative review stage.

\textbf{Iowa-versus-bisect comparison}: \$3,000,000 versus \$52,000. The algorithmic approach costs approximately 1.7\% of the Iowa model.

\subsection{The California Model: \$35 Million+ Per Cycle}

California's Citizens Redistricting Commission spent approximately \$35 million on the 2021 redistricting cycle, including:

\begin{itemize}
  \item Commissioner compensation and benefits (14 commissioners, 18 months): approximately \$5.3 million
  \item CCRC staff (executive director, legal counsel, GIS staff, administrative): approximately \$8.5 million
  \item Public engagement (60+ hearings, translation services, website, outreach): approximately \$6.2 million
  \item Mapping contractors and technology: approximately \$3.1 million
  \item Legal costs (constitutional review, VRA analysis, litigation defense): approximately \$7.5 million
  \item Expert witnesses and demographic analysis: approximately \$2.8 million
  \item Translation and accessibility compliance: approximately \$1.6 million
\end{itemize}

California's high cost reflects its comprehensive process---60+ public hearings, extensive translation services, full-time commissioners for 18 months---and its constitutional entrenchment requirements. The CCRC's maps were challenged in litigation, adding further legal costs.

\textbf{California-versus-bisect comparison}: \$35,000,000 versus \$52,000. The algorithmic approach costs approximately 0.15\% of the California model.

\subsection{Redistricting Litigation: The Hidden Cost}

The cost comparison above excludes redistricting litigation costs, which are the dominant expense in states where maps are challenged. Table~\ref{tab:litigation-costs} summarizes publicly reported litigation costs in major redistricting cases since 2010.

\begin{table}[h]
\centering
\caption{Redistricting litigation costs in selected states, 2010--2024}
\label{tab:litigation-costs}
\begin{tabular}{@{}lrl@{}}
\toprule
State & Cost & Notes \\
\midrule
North Carolina & \$37M+ & Multiple cycles; LWV, NAACP cases \\
Pennsylvania   & \$18M+ & \textit{LWV v. PA}; Senate and House maps \\
Alabama        & \$22M+ & Two cycles; \textit{Milligan} and predecessors \\
Maryland       & \$12M+ & 6th district litigation \\
New York       & \$15M+ & \textit{Harkenrider}; multiple maps \\
Ohio           & \$25M+ & Multiple rounds of maps; seven Supreme Court orders \\
Wisconsin      & \$11M+ & Multiple cycles; legislative and congressional \\
Texas          & \$35M+ & DOJ, private litigation; three census cycles \\
\midrule
\textbf{8-state total} & \textbf{\$175M+} & \\
\bottomrule
\end{tabular}
\end{table}

No state using algorithmic redistricting with a verifiable manifest has faced redistricting litigation over map authenticity. The \texttt{bisect} manifest's byte-reproducibility standard eliminates the ``is this map neutral?'' disputes that generate most redistricting litigation costs.

\subsection{Cost-Benefit Summary}

The cost-benefit case for algorithmic redistricting is compelling on pure financial grounds:

\begin{itemize}
  \item Algorithmic redistricting (\$50K) versus Iowa model (\$3M): 60$\times$ cost reduction
  \item Algorithmic redistricting (\$50K) versus California model (\$35M): 700$\times$ cost reduction
  \item Algorithmic redistricting (\$50K) versus redistricting litigation (\$10M--\$40M): 200--800$\times$ cost reduction
\end{itemize}

Even if implementation costs are three times higher than estimated (due to state-specific complications, higher legal review needs, or first-cycle inefficiencies), the cost comparison is overwhelming. States that have spent double-digit millions on redistricting litigation could have funded algorithmic redistricting for 100 cycles with the same budget.
